\documentclass[11pt]{article}
\usepackage[margin=1.1in]{geometry}
\usepackage{amsmath,amssymb,booktabs}
\usepackage{algorithm}
\usepackage{algpseudocode}
\usepackage[hidelinks]{hyperref}

\newcommand{\Ftwo}{\mathbb{F}_2}
\newcommand{\Feight}{\mathbb{F}_{2^8}}
\newcommand{\Fbig}{\mathbb{F}_{2^{128}}}
\newcommand{\eq}{\mathrm{eq}}
\newcommand{\code}[1]{\texttt{#1}}

\title{Verified Optimizations to Flock's RS Zerocheck Prover\\
{\large measurements on Apple M1 Max, single-threaded, $2^{18}$ BLAKE3 compressions}}
\author{branch \code{snark-fast-improvements}}
\date{August 2026}

\begin{document}
\maketitle

\begin{abstract}
Six changes to the RS (additive-NTT) zerocheck path, each kept only after a
paired, order-alternating A/B measurement with a decisive sign test. Together
they reduce zerocheck round~1 from $1477\,$ms to $1205\,$ms ($-18.4\%$),
round~2 from $433$ to $312\,$ms ($-28\%$), and whole-proof time by
$\approx9\%$ single-threaded. Three are ideas ported
from an independently optimized fork of this prover (the ``challenge'' tree),
re-derived and re-implemented rather than copied; three originated here. The
recurring themes: reductions in fields of characteristic two are
$\Ftwo$-linear and can be deferred almost arbitrarily; work that a circuit
fixes structurally can be skipped with a one-word guard; and on Apple silicon
the boundary between the vector and general register files costs more than
most instruction counts.
\end{abstract}

\section{Setting and notation}

Fields. $\Feight = \Ftwo[x]/(x^8+x^4+x^3+x+1)$ and the GHASH field
$\Fbig = \Ftwo[X]/(X^{128}+X^7+X^2+X+1)$, with the ring embedding
$\varphi:\Feight\hookrightarrow\Fbig$. Write $\alpha := \varphi(x)$ and
$\gamma := X$.

The witness is $z\in\Ftwo^{2^m}$; at the benchmark shape $m=32$
($2^{18}$ BLAKE3 compressions, $k_{\log}=14$ bits per block). After the
univariate skip, round~1 of the zerocheck views the $m$ index bits as
\[
[\;\underbrace{\lambda}_{6\ \text{(univariate)}}\;|\;
 \underbrace{s}_{3\ \text{(small)}}\;|\;
 \underbrace{b}_{4\ \text{(medium)}}\;|\;
 \underbrace{t}_{m-13\ \text{(rest)}}\;],
\]
and must emit, for each of the $64$ points $\lambda$ of the skip domain, the
quadratic message $P^{AB}(\lambda)$ and the linear message $P^{C}(\lambda)$,
plus a pair of \emph{banks} (the ring-switch helper $\hat s_{v,C}$) that keep
the small parity dimension unfolded.

The protocol pins the small and medium challenges to \emph{friendly} values:
$r_{6+i} = \varphi(v_i)$ with $v=(\mathtt{0xF7},\mathtt{0x53},\mathtt{0xB5})$,
and $\beta_i = \gamma^{2^{i-1}}/(1+\gamma^{2^{i-1}})$. These make the eq
weights geometric:
\begin{equation}
\prod_{i=0}^{2}\eq(r_{6+i},K_i) \;=\; C_s\,\alpha^{K},
\qquad
\prod_{i=1}^{4}\eq(\beta_i,b_i) \;=\; \gamma^{b}/D,
\label{eq:friendly}
\end{equation}
with $K=\sum K_i 2^i$, $C_s=\prod_i(1+r_{6+i})=\varphi(\mathtt{0x1C})$, and
$D=\prod_i(1+\gamma^{2^{i-1}})$. The baseline kernels exploit
\eqref{eq:friendly} through lookup tables; several optimizations below exploit
it in the opposite direction, replacing tables by folds at the actual
challenge values.

\paragraph{Measurement protocol.}
All numbers are single-threaded at the shape above. Each verdict is a paired
A/B: both arms built once, one discarded warm-up each, then eight pairs with
the \emph{order alternating within each pair} (a fixed order is biased by
thermal drift by roughly the size of the effects measured here). Phase times
are the minimum over the $\sim$5 zerocheck invocations within one run. A
change is kept only on a decisive sign test; two of the six wins were
$8/8$ with disjoint ranges. Measure only on AC power with the battery above
$\sim60\%$: low battery caps frequencies, and fast-charging a nearly empty
battery is worse.

\section{Deferred reduction in vector registers}\label{sec:wideneon}

For $u\in\Ftwo^{256}$ (an unreduced carry-less product) let
$R(u)\in\Fbig$ be reduction mod $X^{128}+X^7+X^2+X+1$. $R$ is
$\Ftwo$-linear, so for any products $u_i = a_i \odot b_i$ (unreduced),
\[
\textstyle\sum_i R(u_i) \;=\; R\!\left(\sum_i u_i\right),
\]
and a sum of $n$ field products needs one reduction, not $n$. The baseline
already used this identity, but held the $256$-bit accumulator as a four-word
struct in \emph{general} registers: each product paid six lane extracts to
leave the vector file, and each accumulate was four scalar XORs.

The fix is representational, not algorithmic: keep the accumulator as two
$128$-bit vector registers. The unreduced product uses Karatsuba,
\[
(a_l + a_h X^{64})(b_l + b_h X^{64}) : \quad
\mathit{ll}=a_l b_l,\ \ \mathit{hh}=a_h b_h,\ \
\mathit{cross}=(a_l{+}a_h)(b_l{+}b_h)+\mathit{ll}+\mathit{hh},
\]
three \code{PMULL}s instead of the schoolbook four; accumulation is two vector
XORs; the four extracts are paid once per worker chunk on the way out, and the
existing scalar reducer is reused unchanged.

\emph{Measured: round-2 phase $433.1\to375.2\,$ms ($-13.4\%$).}

\section{Structurally fixed rows of the $b$ operand}\label{sec:zeroskip}

Round~1's dominant kernel transforms the operands $a,b$ through a table $T$
implementing the inverse-NTT composition; $T$ is $\Ftwo$-linear, so
$T(0)=0$. A census of the packed BLAKE3 witness (256 word positions per
block $\times$ 256 blocks $\times$ 3 independent witnesses) shows the circuit
pins $38$ of $256$ eight-byte $K$-rows of $b$ regardless of the inputs, taking
only three values:
\[
\underbrace{\mathtt{0xff}^8}_{22\ \text{positions (const-one wires)}},\qquad
\underbrace{0^8}_{15\ \text{positions (structural zeros)}},\qquad
\text{one mixed word.}
\]
For a zero row the entire row's work vanishes: $db=T(0)=0$, hence
$y = da \odot db = 0$ contributes nothing to any accumulator. The kernel
gains a single guard:

\begin{algorithmic}[1]
\State \textbf{if} $\mathrm{load}_{64}(b\_\mathit{row}) = 0$ \textbf{then
return} \Comment{skips all 64 table loads and 4 $\Feight$ multiplies}
\end{algorithmic}

\noindent
No census data ships and no positions are tracked: the guard is exact for any
witness, and a disagreeing witness falls through to the generic path.
(The all-ones rows were also tried: constant-folding them saves only the $b$
half of a row's work, and halving a row's loads halves its memory-level
parallelism, returning $6\,$ms where $27$ were predicted. Reverted; whole-row
elimination is what pays.)

\emph{Measured: round~1 $-42.4\,$ms ($-2.9\%$), $8/8$.}

\section{Instruction-level parallelism in the gather drain}\label{sec:twolane}

The drain accumulates, per output lane, three XOR chains of depth
$n_{\mathrm{med}}=16$ (one per bank), each link a table gather feeding an XOR.
The gathers are independent but the chains are serial, so one lane exposes
three dependency chains to the out-of-order engine. Processing two lanes per
iteration doubles that to six with no change in total work:

\begin{algorithmic}[1]
\For{$\ell = 0,2,4,\dots,62$}
  \State $u_0,u_1,u_2 \gets 0$;\quad $v_0,v_1,v_2 \gets 0$
    \Comment{banks for lanes $\ell$ and $\ell{+}1$}
  \For{$b = 0 \dots n_{\mathrm{med}}-1$}
    \State $u_j \mathrel{\oplus}= \mathrm{convert}[b][\cdot]$;\quad
           $v_j \mathrel{\oplus}= \mathrm{convert}[b][\cdot]$
           \Comment{6 independent chains in flight}
  \EndFor
\EndFor
\end{algorithmic}

This targets latency, not throughput: an attempt to shrink the gather
\emph{table} ($64\to8\,$KiB, doubling gather count) had cost $3.6\%$ --- the
table already fit M1's $128\,$KiB L1D, establishing that the drain is bound by
gather count and chain depth, not footprint.

\emph{Measured: round~1 $-63.1\,$ms ($-4.3\%$), $8/8$. Combined with
Section~\ref{sec:zeroskip}: $1477.3\to1403.1\,$ms ($-5.0\%$).}

\section{Unreduced accumulation with multiplicative weight splitting}
\label{sec:pmull}

The prep kernel computes, per $K$-row ($K\in\{0,\dots,7\}$) and 16-lane block,
$y_K = a_K \odot b_K \in \Feight$ and accumulates $\sum_K x^K y_K$ in 16-bit
lanes, reducing once at the end. The baseline computed each $y_K$
\emph{reduced}: two \code{PMULL}s for the raw product plus a reduction costing
four more, then a widening shift by $K$ --- six \code{PMULL}s per row-block for
a reduction that the final fold makes redundant.

Let $u_K$ be the \emph{unreduced} 15-bit product. Then
$\sum_K x^K u_K \equiv \sum_K x^K y_K \pmod{p_8}$, but $x^K u_K$ overflows 16
bits for $K\ge2$. Split the weight so every factor lands where it is free:
\[
x^K \;=\;
\underbrace{(x^4)^{[K\ge4]}}_{\text{choice of gather table}}\cdot
\underbrace{(x^2)^{[\lfloor K/2\rfloor\ \mathrm{odd}]}}_{\text{byte-wise
doubling of the reduced operand}}\cdot
\underbrace{x^{K\bmod 2}}_{\text{in-lane shift}}.
\]
The $x^4$ factor costs nothing per element: a second table image with every
byte pre-multiplied by $x^4$ scales each gathered XOR-sum by $x^4$, by
linearity of $T$. The $x^2$ factor is a six-instruction shift-and-conditional-XOR
on the already-reduced 8-bit operand ($\deg$ still $\le7$). The residual shift
is one vector instruction and raises term degree to at most $15$.

\begin{algorithm}[h]
\caption{Row-block accumulate, before and after}
\begin{algorithmic}[1]
\State \textbf{before:}\ $\mathit{acc} \mathrel{\oplus}=
  \mathrm{shift}_K\!\big(\mathrm{reduce}(\mathrm{pmull}(da, db))\big)$
  \Comment{$6$ \code{PMULL}}
\State \textbf{after:}\ \ $da \gets [K\ge4]\ ?\ \text{gathered from }x^4\text{
  table} : da$;\quad $da \gets [(K/2)\bmod 2]\ ?\ x^2\!\cdot\!da : da$
\State \phantom{\textbf{after:}}\ $\mathit{acc} \mathrel{\oplus}=
  \mathrm{shift}_{K\bmod2}\!\big(\mathrm{pmull}(da, db)\big)$
  \Comment{$2$ \code{PMULL}}
\end{algorithmic}
\end{algorithm}

Correctness needs the final reducer to be exact to degree $15$, one beyond its
documented ``$\le14$''; both the scalar and vector reducers were verified
exhaustively over all $2^{16}$ inputs (tests now permanent). Gather traffic is
unchanged apart from the second $16\,$KiB table image.

\emph{Measured: round~1 $1401.4\to1273.4\,$ms ($-9.1\%$), $8/8$, every pair in
$[-8.7\%,-10.1\%]$. The reference implementation's attribution predicted
$100$--$140\,$ms; measured $128$.}

\section{The C-side message as a fold of the lincheck stripe}
\label{sec:stripe}

In the fast hash setups $C=I$, so $\hat c = \hat z$ and everything round~1
emits on the C side is \emph{linear} in the witness. Concretely the two banks
are, for parity $p\in\{0,1\}$ and lane $\lambda$,
\begin{equation}
\mathrm{bank}_p[\lambda] \;=\;
\sum_{K\equiv p\ (2)} \alpha^{K}
\sum_{b} \frac{\gamma^{b}}{D}
\sum_{t} \eq(r_{\ge13}, t)\; z[\lambda,K,b,t],
\label{eq:banks}
\end{equation}
a multilinear evaluation --- which the baseline nevertheless computed with the
quadratic side's machinery: a bit transpose of every $8192$-bit window plus
$32$ of the drain's $48$ gathers per lane.

By \eqref{eq:friendly}, a \emph{true} multilinear fold at the actual pinned
challenges reproduces the weights in \eqref{eq:banks} exactly. The prover
already materializes the \emph{lincheck stripe}, a repacking of $z$ indexed
$[\,i_{\mathrm{inner}}\,|\,i_{\mathrm{outer}}\,]$, for the later lincheck
phase. So:

\begin{algorithmic}[1]
\State $v \gets \mathrm{fold\_stripe\_outer}(z_{\mathrm{stripe}},\
  \eq(r_{k_{\log}..m}))$
  \Comment{the one $O(2^m)$ pass; same kernel lincheck uses}
\For{$d = k_{\log}-1$ \textbf{downto} $7$}
  \Comment{data halves each round: $2^{14}\to2^{7}$ elements}
  \State $v[i] \gets v[i] + r_d\,(v[i]+v[i+2^{d}])$
\EndFor
\State $\mathrm{bank}_p[\lambda] \gets \dfrac{\eq(r_6,p)}{C_s}\;
       v[\,p\cdot64+\lambda\,]$
  \Comment{$C_s=(1{+}r_6)(1{+}r_7)(1{+}r_8)$, from $r$ itself}
\end{algorithmic}

\noindent
Dimension 6 (the parity) and dimensions $0..5$ (the lanes) are kept unfolded;
the constant in step~3 undoes the two eq factors the partial fold applied,
derived from the identity $Y_p = (C_s/\eq(r_6,p))\,\mathrm{bank}_p$ obtained
by factoring \eqref{eq:friendly} over the folded dimensions. With the banks
gone from the drain, the workers run AB-only: the per-window bit transpose is
deleted and the drain issues one gather per (lane, $b$) instead of three.
Equivalence was proven at two levels before wiring: banks bit-identical to the
drain's, and all three entry outputs identical, dense and padded.

\emph{Measured: round~1 $1277.5\to1204.7\,$ms ($-5.7\%$), $8/8$. The cost of
the added fold ($\approx180\,$ms) is included; the removals
($\approx250\,$ms) net out.}

\section{Register-resident round 2}\label{sec:qres}

Round~2's message loop crossed the vector/general register-file boundary three
times per output pair: the table-driven fold extracted its vector accumulator
into a two-word struct; the message products
$g_1 = a_1 b_1$, $g_\infty = (a_0{+}a_1)(b_0{+}b_1)$ went through the scalar
field multiply (struct in, struct out); and the accumulate of
Section~\ref{sec:wideneon} re-entered the vector file. The rewrite keeps the
whole chain in vector registers: the fold returns its register unextracted;
the products use an in-register Karatsuba multiply (five \code{PMULL}s ---
three for the product, two for the fold through $X^{128}\equiv X^7{+}X^2{+}X{+}1$
--- against the baseline's six); the $\eq\cdot g$ products feed the
Section~\ref{sec:wideneon} accumulator directly. Per pair, the only remaining
memory traffic is the four required stores of folded values and one eq load.

\emph{Measured: round-2 phase $358.1\to311.6\,$ms ($-13.0\%$), $8/8$, every
pair in $[-12.6\%,-13.3\%]$ --- the tightest distribution of the whole
effort, taken on a machine with an active browser session (the min-of-5
protocol absorbs interactive bursts).}

\section{Summary}

\begin{table}[h]
\centering
\begin{tabular}{@{}llrr@{}}
\toprule
Optimization & Kind & Phase effect (ST) & Sign test\\
\midrule
\S\ref{sec:wideneon}\ deferred reduction in registers & representation &
  round 2: $-13.4\%$ & decisive\\
\S\ref{sec:zeroskip}\ structurally-zero $b$ rows & work removal &
  round 1: $-2.9\%$ & $8/8$\\
\S\ref{sec:twolane}\ two-lane drain & ILP &
  round 1: $-4.3\%$ & $8/8$\\
\S\ref{sec:pmull}\ unreduced accumulate, split weights & arithmetic &
  round 1: $-9.1\%$ & $8/8$\\
\S\ref{sec:stripe}\ C side as stripe fold & structural &
  round 1: $-5.7\%$ & $8/8$\\
\S\ref{sec:qres}\ register-resident round 2 & representation &
  round 2: $-13.0\%$ & $8/8$\\
\midrule
\multicolumn{2}{@{}l}{round 1 cumulative} &
  \multicolumn{2}{r@{}}{$1477 \to 1205\,$ms\quad($-18.4\%$)}\\
\multicolumn{2}{@{}l}{round 2 cumulative
  (\S\ref{sec:wideneon}+\S\ref{sec:qres})} &
  \multicolumn{2}{r@{}}{$433 \to 312\,$ms\quad($-28\%$)}\\
\multicolumn{2}{@{}l}{end-to-end, all six} &
  \multicolumn{2}{r@{}}{$52{,}110 \to 56{,}820$ comp/s\quad($\approx+9\%$)}\\
\bottomrule
\end{tabular}
\end{table}

\paragraph{Negative space.}
The kept changes are one side of a 17-experiment record; the failures shaped
them. On this core, re-encoding fixed work never paid: pairing multiplies
through a memory-interfaced kernel ($+10\%$), a four-register structured load
($+12\%$, microcoded), three-way XOR fusion ($+1.8\%$), an $8\times$ smaller
gather table ($-3.6\%$ regression), and deferring an already-cheap reduction
($0\%$) all failed, while every win either removed work, added independent
work in flight, or removed register-file boundary crossings. Two further
structural transplants measured neutral end-to-end and were reverted despite
making the zerocheck \emph{bucket} look better (hoisting the challenge-independent
prep under the commit; the lookahead double-round, which loses $7$--$15\%$ in
the zerocheck tail on both this machine and the reference tree's, while
\emph{helping} the PCS open by $7\%$ --- the same technique, opposite signs,
two call sites).

\end{document}
